\chapter{Einleitung}
Ein Debugger ist ein grundlegendes Werkzeug zur Entwicklung komplexer Systeme und Programme. Er ermöglicht es unter anderem, Programme anzuhalten und Speicher auszulesen. Für Anwendungen, die auf etablierten Betriebssystemen laufen, stehen Debuggingwerkzeuge, wie der GDB (GNU Debugger), zur Verfügung. Es kann aber sein, dass auf einem System kein vollständiger Debugger ausgeführt werden kann. Dazu zählen beispielsweise Betriebssystemkernel oder Betriebssysteme, die noch nicht alle Funktionen unterstützen, die ein Debugger benötigt. In solchen Fällen kann Remote-Debugging verwendet werden, um dennoch das System auf Fehler untersuchen zu können \cite{gdb}.

\section{Motivation}
D3OS ist ein Betriebssystem, das von der Abteilung für Betriebssysteme der HHU entwickelt wird. Es kann aktuell mithilfe von QEMU und GDB debuggt werden \cite{d3os-repo}. QEMU ist ein Open-Source Systememulator, beziehungsweise Virtualisierungsprogramm. Es stellt eine Debugging-Schnittstelle bereit, über die GDB zum Beispiel den emulierten Systemzustand kontrollieren, Speicher auslesen und Breakpoints setzen kann \cite{qemu}. Dadurch kann D3OS in einer emulierten Umgebung auf Fehler untersucht werden.\\
Die Ausführung in einem Emulator unterscheidet sich jedoch von der Ausführung auf echter Hardware. QEMU abstrahiert und emuliert Hardwareverhalten. Es kann aber sein, dass bestimmte Fehler von zum Beispiel Hardwareeigenschaften oder Timings abhängig sind. Manche Fehler könnten also durch die emulierte Hardware verhindert werden, aber auf echter Hardware auftreten. In solchen Fällen reicht der QEMU-Debugger nicht.\\
D3OS soll nicht nur in QEMU, sondern auch auf realer Hardware auf Fehler untersucht werden können. Da auf dem Zielsystem kein vollständiger Debugger ausgeführt werden kann, muss der Kernel eine Schnittstelle bereitstellen, über die GDB mit ihm kommunizieren kann. Das kann durch einen GDB-Remote-Stub erreicht werden, der als Kernel-Thread von D3OS läuft und dadurch Remote-Debugging ermöglicht.

\section{Zielsetzung}
Ziel dieser Arbeit ist die Entwicklung und Integration eines GDB-Stubs für D3OS. Dieser soll es ermöglichen, das Betriebssystem und darauf ausgeführten Code auf echter Hardware zu debuggen. Der Stub kommuniziert über das GDB Remote Serial Protocol mit einer GDB Instanz, die auf einem anderen Host-Rechner ausgeführt wird. Als Grundlage für die Kommunikation wird die Rust Crate \texttt{gdbstub} verwendet, die von GDB empfangene Pakete verarbeitet und entsprechende Funktionen aufruft \cite{gdbstub}. Die Pakete werden über den seriellen Port ausgetauscht.\\
Der Stub soll im Kernel laufen und grundlegende Debugging-Funktionen bereitstellen. Dazu gehören das Lesen und Schreiben von Registern, Zugriff auf Speicher, das Setzen und Entfernen von Breakpoints sowie die Unterstützung von Single-Stepping. 

\section{Verwendete Tools und Entwicklungsumgebung}
Der Stub wurde in einer Kopie des Repository von D3OS entwickelt \cite{d3os-repo}, welches ebenfalls stetig weiterentwickelt wurde. Während der Entwicklung des Stubs wurde der Scheduler von D3OS von einem Single-Core-Scheduler zu einem Multi-Core-Scheduler geändert. Da der Stub die Kontrolle über alle laufenden Threads braucht, ist er auf Funktionen des Schedulers angewiesen, die sich mit der Umstellung anders verhalten. Auf der neusten Version des Stubs funktioniert dieser mit dem neuen Scheduler in QEMU, aber nicht auf echter Hardware. Die Probleme die dabei auftreten werden im Kapitel Evaluation beschrieben. Die nutzbare Version des Stubs basiert auf einer älteren Version von D3OS.\\
Als Entwicklungsumgebung wurde Visual Studio Code \cite{vscode} ohne Add-Ons verwendet. Zum testen der graphischen Nutzeroberfläche von VS Code wurden die entsprechenden Extensions zur Nutzung von GDB in VS Code verwendet. Um möglichst schnell und effizient neue Funktionen zu testen, wurde QEMU verwendet. Der Stub wurde aber auch auf echter Hardware auf den Bildungsrechnern der Universität getestet.\\
ChatGPT \cite{chatgpt} wurde zur Fehlersuche und als Suchmaschine während dieser Arbeit verwendet.